%!TEX root = ../template.tex
%% Capitulo 8 -- Conclusoes e trabalho futuro. Escrito em 2026-09-17.
%% Numeros: Caps. 3, 5, 6 e 7 (nao ha numeros novos neste capitulo).
%% Versao anterior (esqueleto) em _revisao/2026-09-17-cap8-antes.tex.

\chapter{Conclusões e trabalho futuro}
\label{cha:conclusoes}

A dissertação abriu com uma pergunta simples de fazer e difícil de
responder: consome a unidade hoje mais do que a referência prevê para a carga
que tem? No início, a resposta estava espalhada por três sistemas e chegava
tarde. No fim, a plataforma dá o número todos os dias, reconciliado e
rastreável até à fonte. O que a dissertação acrescenta é o resto da resposta:
esse número só significa alguma coisa se se souber em que domínio a
referência vale, contra que comparador foi medida, com que cobertura as
bandas funcionam e desde quando está em vigor. Na unidade estudada, a maior
parte dessa informação não existia, e onde foi medida mostrou que a
referência em produção não serve da mesma maneira aos cinco vetores.

\section{Conclusões}
\label{sec:conclusoes}

\paragraph{Construir.} As três fontes de dados da refinaria --- historiador
de processo, balanços mensais de produção e de utilidades, relatórios de
eletricidade --- ficaram reunidas num modelo dimensional único, que cobre
dezasseis unidades e é atualizado diariamente. A validação a três leituras,
em 400 células unidade--vetor--mês de janeiro a maio de 2026, explicou as 94
divergências iniciais por sete causas, com a última correção na versão~15.1
do mapeamento; a atualização seguinte correu sem erros e sem divergências
detetadas. As causas estavam nas três fontes: no
motor de ingestão, no mapeamento e no próprio relatório oficial. O fuel gás
ficou validado apenas contra duas fontes que leem o mesmo balanço, o que é
consistência e não verificação (Secção~\ref{sec:validacao-dados}). A
primeira camada da fiabilidade, a dos dados, ficou resolvida no âmbito
validado, que são cinco meses de 2026: nos seis anos avaliados, o canal de
\emph{tags} e o de balanço afastam-se de forma variável de ano para ano, e a
série avaliada é a do balanço (Secção~\ref{sec:modelo-dimensional}). A lição
transferível é que nenhuma das três fontes, isolada, era fiável.

\paragraph{Diagnosticar.} O método de referência em produção foi replicado
de forma exata nas 30 células vetor--ano de 2020 a 2025, com erro máximo da
ordem de $10^{-10}$ contra uma tolerância de $10^{-6}$
(Secção~\ref{sec:replicacao}). Avaliado depois contra um protocolo fixado
antes dos resultados, o método mostra quatro comportamentos diferentes nos
quatro vetores diários. O fuel gás tem conteúdo preditivo: aplicado ao ano
seguinte, reduz o erro absoluto médio em 30 a 40\% face à média do ano de
treino em quatro de cinco transições. Mas o declive desce cerca de 30\% em
cinco anos, e a banda de alarme perde cobertura nos pares que atravessam
essa deriva. Nos três níveis de vapor, a carga não traz informação
utilizável de um ano para o seguinte: o \emph{skill} do vapor de 24 e de
10~bar fica próximo de zero, e o do vapor de 3~bar muda de sinal. O vapor de
24~bar passa os testes de estabilidade e de calibração precisamente por não
haver relação a mudar; o de 10~bar tem um nível que sobe de ano para ano; o
de 3~bar tem a pior calibração, com coberturas entre 0,38 e 0,997. A
eletricidade é uma alocação mensal repartida pelos dias, e o alarme diário
sobre ela não tem resolução para ver um desvio intramensal. Três observações
atravessam os vetores: o limiar de carga de $18\,600$~t/d exclui 10,4\% dos
dias, que pertencem a outro regime operatório e definem um domínio de
validade nunca declarado; os resíduos têm autocorrelação de primeira ordem
entre 0,46 e 0,90, o que invalida os erros-padrão clássicos em que a banda
se apoia; e a reestimação anual acomoda nos parâmetros a deriva que o
indicador devia mostrar. As covariáveis de processo disponíveis (densidade e condições de
\emph{flash}) não são um remédio geral: acrescentam valor em quatro de cinco
pares no vapor de 10~bar, com um ganho frágil à análise de sensibilidade,
oscilam entre ganho robusto e prejuízo no de 3~bar, e prejudicam na maioria
dos pares no fuel gás e no de 24~bar. A detetabilidade real ficou por estimar, porque o
calendário operacional não estava disponível. Estes resultados são
moderadamente robustos às convenções: mudar o limiar em 20\% altera 20 de
208 rótulos, quase todos em pares fronteiriços
(Secção~\ref{sec:sintese-modos-falha}).

\paragraph{Prescrever.} As linhas orientadoras levam a classe linear
estática tão longe quanto a demonstração permitiu: domínio declarado nas variáveis explicativas,
exclusões por causa e nunca por valor, especificação antes de robustificação,
estimação robusta pré-especificada, bandas que admitem três desfechos
(incluindo nenhuma banda suportada), critério de aceitação encadeado no
tempo e recalibração por evidência e não por calendário
(Secção~\ref{sec:guidelines}). Aplicado aos 32 candidatos avaliados, o
critério de aceitação fixado antes dos resultados não aceitou nenhum: o fuel
gás cai pela banda, os vapores pelo conteúdo preditivo. O resultado depende
da regra de quatro pares em cinco, que estava registada e não foi alterada;
com três em cinco, três dos 16 candidatos de mínimos quadrados passariam.
Não demonstra que a classe seja impossível; demonstra um procedimento que
publica um resultado negativo em vez de o ajustar. A simulação com verdade
conhecida acrescenta uma conclusão com consequência prática: sob o gerador
usado, trocar o monitor de nível por uma soma acumulada aumenta a deteção de
um degrau de uma escala em 32 a 49 pontos percentuais, enquanto trocar a
referência linear por um modelo aditivo muda-a em menos de três. Com a
autocorrelação que os vapores exibem, porém, um degrau desses é detetado
dentro de noventa dias em cerca de metade das réplicas e uma rampa num
quarto (Secção~\ref{sec:verdade-conhecida}). A especificação que volta à
plataforma --- duas pistas, domínio visível, calibração monitorizada, ficha
por referência e políticas separadas --- está proposta e não implantada
(Secção~\ref{sec:implicacoes-plataforma}).

\paragraph{A revisão.} Nas 331 publicações elegíveis, o rácio de
intensidade é a família de modelo mais frequente, embora a regressão apareça
em quase metade das que têm classificação concordante; a ISO~50001 é citada
em 185 de 258 e a ISO~50006 em 48. Nos 39 estudos de processo contínuo dos
setores pré-registados, a energia raramente é descrita como medida na
fronteira do modelo, o domínio de validade quase nunca é declarado e a
validação fora da amostra aparece em quatro. Os estudos relatam o efeito
sobre o modelo e quase nunca a causa física no processo
(Secção~\ref{sec:gap}). Nenhum dos 13 estudos com aplicação em refinaria ou
petroquímica junta medição na fronteira, regra de reestimação, domínio
declarado e teste fora da amostra. O caso de Sines não é, portanto, uma
exceção: é uma instância documentada do que a literatura deixa em aberto.

\paragraph{O que isto diz sobre fiabilidade.} O Capítulo~\ref{cha:introducao}
definiu uma plataforma fiável por quatro condições verificáveis. A
Tabela~\ref{tab:fiabilidade-estado} resume onde cada uma ficou. A primeira
fica cumprida pela plataforma; a terceira fica cumprida pelo diagnóstico, que
não existia antes; a segunda e a quarta ficam especificadas, e é aí que está
o trabalho que falta. As normas de referência exigem indicadores,
referências, normalização e resposta a desvios, e deixam à organização a
demonstração de que a referência merece a comparação. A fiabilidade exige as
duas coisas: dados reconciliados e referências avaliadas.

\begin{table}[htbp]
  \centering\small
  \caption[Condições de fiabilidade: estado no fim do trabalho]{As quatro
  condições de fiabilidade definidas no Capítulo~\ref{cha:introducao} e o seu
  estado na unidade estudada no fim do trabalho.}
  \label{tab:fiabilidade-estado}
  \begin{tabularx}{\textwidth}{@{}>{\raggedright\arraybackslash}p{3.3cm}>{\raggedright\arraybackslash}X>{\raggedright\arraybackslash}X@{}}
    \toprule
    Condição & Antes da plataforma & No fim do trabalho \\
    \midrule
    Cada número é rastreável até à fonte e à regra & Folhas de cálculo manuais; três fontes separadas & Cumprida no âmbito validado (janeiro a maio de 2026): modelo único, mapeamento como contrato, validação a três leituras, proveniência por \emph{hash} \\
    \addlinespace
    O domínio da comparação está declarado & Implícito no limiar de carga & Medido (10,4\% dos dias fora do domínio, outro regime); declaração e supressão do alarme especificadas, não implantadas \\
    \addlinespace
    O desempenho da referência está quantificado & $R^2$ de treino & Cumprida para 2020--2025: \emph{skill} contra comparadores triviais, estabilidade, calibração e valor das covariáveis por vetor \\
    \addlinespace
    A incerteza acompanha o resultado & Banda $\pm2\sigma$ assumida & Cobertura medida (0,38 a 0,997); bandas com três desfechos e cobertura monitorizada especificadas, não implantadas; detetabilidade real não estimada \\
    \bottomrule
  \end{tabularx}
\end{table}

As limitações que condicionam estas conclusões estão na
Secção~\ref{sec:limitacoes} e não se repetem aqui. Duas pesam mais do que as
outras sobre o que se pode fazer a seguir: sem calendário operacional nenhum
desvio pode ser classificado e nenhuma taxa de deteção real estimada; e a
decisão de aceitação assenta em cinco pares anuais dependentes, cuja
potência não foi calculada.

\section{Roteiro e trabalho futuro}
\label{sec:roteiro-modelos}
\label{sec:trabalho-futuro}

Os modelos que responderiam à não estacionariedade --- com estado de
degradação, filtros de nível e declive, modelos aditivos com covariáveis de
processo, modelos hierárquicos entre unidades, modelos físicos da bateria de
pré-aquecimento --- não formam uma escada em que o mais complexo é melhor. A
simulação mostrou que, sob o gerador usado, a escolha do monitor pesa mais do
que a da referência, e o critério de aceitação mostrou que um modelo mais
complexo sem as condições de avaliação não seria mais aceitável do que o
atual. O roteiro organiza-se por isso a partir das falhas observadas, e começa
pelas condições que tornam qualquer modelo avaliável
(Tabela~\ref{tab:roteiro}).

% 2026-09-20: com a linha nova sobre a fronteira do indicador a tabela
% deixou de caber numa pagina e transbordava. Passa a xltabular, como as
% Tabelas 1.1, 2.1 e 2.4--2.6: parte-se com cabecalho repetido.
{\footnotesize
\begin{xltabular}{\textwidth}{@{}>{\raggedright\arraybackslash}p{2.6cm}>{\raggedright\arraybackslash}X>{\raggedright\arraybackslash}p{2.6cm}>{\raggedright\arraybackslash}p{2.4cm}>{\raggedright\arraybackslash}X@{}}
\caption[Roteiro a partir das falhas observadas]{Roteiro de trabalho futuro organizado pelas falhas observadas. As três primeiras linhas são pré-condições das seguintes.}\label{tab:roteiro}\\
\toprule
Falha observada & Opção & Dados adicionais & Comparador & Critério de sucesso \\
\midrule
\endfirsthead
\toprule
Falha observada & Opção & Dados adicionais & Comparador & Critério de sucesso \\
\midrule
\endhead
\midrule
\multicolumn{5}{r}{\emph{continua na página seguinte}}\\
\endfoot
\bottomrule
\endlastfoot
    Detetabilidade real não estimável; nenhum desvio classificável & Integrar o calendário operacional e rótulos de eventos na plataforma & Paragens, arranques, limpezas, mudanças de crude, intervenções de instrumentação & Situação atual, sem eventos & Taxa de deteção e atraso medidos sobre eventos reais; alarmes classificados \\
    \addlinespace
    Critério de aceitação sem margens materiais & Elicitar com a operação a menor mudança que importa e a precisão mínima & Margens por vetor; custo de um falso alarme e de um desvio não visto & Regra de 4 em 5 atual & Critério de equivalência ou não inferioridade fixado antes de avaliar qualquer candidato novo \\
    \addlinespace
    Fronteira e resolução: vapor de 10~bar líquido, eletricidade mensal, definição de carga pendente & Separar produção e consumo de vapor; medir ou submedir a eletricidade diária; fechar a definição de carga & Medidores de produção e consumo por nível; potência elétrica por unidade & Série atual & Séries diárias validadas contra o fecho mensal, como na Secção~\ref{sec:validacao-dados} \\
    \addlinespace
    \novo{Fronteira do indicador: o bloco agrega à coluna secções auxiliares cujo consumo não responde à carga de crude} & \novo{Quantificar a fração do consumo do bloco que vem das secções auxiliares, separando os canais que o mapeamento já distingue; avaliar uma variável de atividade própria para essas secções, ou uma referência por secção} & \novo{Canais já separados no mapeamento mestre; variável de atividade das secções auxiliares (produção de gases, carga de enxofre)} & \novo{Reta única sobre a carga de crude do bloco} & \novo{Fração estimada com intervalo; um desenho alternativo só entra se reduzir a variância residual fora da amostra sem degradar a banda} \\
    \addlinespace
    Fuel gás: conteúdo preditivo com deriva do declive & Piloto prospetivo de referência congelada com soma acumulada, contra um modelo com estado (nível e declive variáveis) e contra uma regressão com um indicador de incrustação calculado da bateria & Temperaturas e caudais dos permutadores; datas de limpeza & Referência anual em produção; referência congelada de 2020 & Passar o critério revisto num ano não usado; não se promete ganho \\
    \addlinespace
    Vapores de 24 e 10~bar: sem relação utilizável com a carga & Tratar uma referência constante bem avaliada como comparador obrigatório; procurar explicação no estado da rede de vapor & Balanço de vapor da refinaria por nível & Média do ano de treino & Um modelo só entra se prever melhor do que a média fora da amostra e mantiver a banda suportada \\
    \addlinespace
    Vapor de 3~bar: relação provavelmente não linear, covariáveis úteis & Modelo aditivo com covariáveis de processo, validado para a frente no tempo & Caudais de vapor de \emph{stripping}, temperaturas de corte & Regressão linear com as mesmas covariáveis; média & Valor incremental com intervalo a excluir zero em anos não usados na afinação \\
    \addlinespace
    Expansão a outras unidades & Modelos hierárquicos com partilha parcial entre unidades semelhantes, depois de o procedimento passar numa unidade & Unidades com catalisador: dias desde a regeneração, temperatura de reação & Modelos independentes por unidade & Mesmo critério de aceitação; análise de custo e benefício da expansão \\ \\
\end{xltabular}
}

Duas linhas de trabalho ficam fora da tabela por não dependerem de modelos.
A primeira é implantar o bloco de saúde das referências na plataforma, com
as duas pistas lado a lado e a validação a três leituras repetida em cada
fecho mensal; é a parte do trabalho que pode entrar em produção sem nova
investigação. A segunda é fechar o que a revisão deixou aberto: verificar por
leitura humana os valores do Corpus~B que sustentam os resultados, fazer o
rastreio de referências que não foi feito e alargar a pesquisa à literatura
de \emph{benchmarking} industrial que as expressões registadas não
apanharam, e confrontar o argumento normativo com a edição de 2023 da
ISO~50006.

A pergunta de abertura pode então ser respondida por completo. A plataforma
diz quanto a unidade consumiu. Dizer se consumiu mais do que devia, à carga
que teve, exige uma referência cujo domínio, desempenho e incerteza estejam
à vista do utilizador. Medir já se faz; é isso que falta para saber.
